\section{Project synthesis}
\label{sec:synthesis}

\begin{synthesis}{where the law of the omitted week actually lives}

\textbf{Claim.} The omission of a week of Ordinary Time binds every book of the Roman Rite, but the
rule that determines \emph{which} week is omitted has its only reasoned statement in the
\latin{Praenotanda} of the \latin{Ordo lectionum Missae}; and the practical weight of the omission
falls very unequally across the books, heavily on the Lectionary and lightly on the Missal, because
the Missal deliberately detached its Ordinary Time weekday formularies from their week numbers and
the Lectionary deliberately did not.

\textbf{Reasons.}

\emph{First, the loci.} The Universal Norms state that Ordinary Time is thirty-three or thirty-four
weeks (\loc{nualc}~43) and where the two runs begin and end (\loc{nualc}~44). They never say how the
weeks are numbered, which one is dropped, or why. The Missal's rubric~2\,b states the omission in a
single subordinate clause --- \latin{omittitur prima hebdomada, quæ sumenda esset post Pentecosten}
--- with no criterion and no reason. Only \loc{olm}~104.3 supplies both: the same clause, followed by
\latin{ut retineantur in fine anni textus eschatologici qui ultimis duabus hebdomadis assignantur},
and then two footnotes working the thirty-four-week and thirty-three-week cases on numbers. The
purpose clause is the only one in this area of law, and it is about texts.

\emph{Second, the asymmetry is drafted, not accidental.} The Missal's rubric~3\,b permits any of the
thirty-four Ordinary Time Masses to be said on any weekday of the season, with the pastoral good of
the faithful in view, and \loc{igmr}~363 widens the choice of orations further. No euchological text
of Ordinary Time is therefore ever made unavailable by an omitted week; at most a priest loses the
default. The Lectionary runs the other way: \loc{olm}~69.4 ties the ferial first reading to a
numbered two-year course; \loc{igmr}~358 directs that the ferial readings be taken on the days to
which they are assigned; and \loc{igmr}~355 tells the priest not to omit them frequently or without
sufficient cause, because the Church desires a richer table of God's word to be spread for the
faithful. One book relaxes the tie between week number and text; the other tightens it. The omission
therefore costs the Missal one Sunday formulary and costs the Lectionary ten weekday sets.

\emph{Third, the reason given fits the asymmetry.} The stated purpose of dropping a week in the
middle rather than at the end is to preserve the eschatological readings of the last two weeks. That
is a concern about the arrangement of Scripture across the year, not about collects. A rule whose
stated criterion is scriptural is most naturally located where Scripture is arranged.

\textbf{The strongest counterargument, at full strength.} The second paragraph of \loc{nualc}~44 says:
\latin{Eadem ratione adhibetur series formulariorum, quæ pro dominicis et feriis huius temporis
invenitur tum in Missali tum in Liturgia Horarum (voll.~III-IV)} --- ``by the same reckoning is used
the series of formularies which for the Sundays and weekdays of this season is found both in the
Missal and in the Liturgy of the Hours, volumes III--IV.'' That sentence is in the Universal Norms,
it is promulgated law, and it says in terms that one and the same reckoning of weeks governs three
books, of which the Lectionary is not even the one named. On this reading the omitted week is a fact
about the \emph{season} --- a missing unit of the temporal cycle --- which every book then records;
the Lectionary is simply the book that happened to explain itself, and treating its explanation as
the seat of the law confuses a commentary with a norm. The counterargument has three further
supports. The Missal's rubric~3\,a does bind the Sunday formulary to the Sunday's number, so the
Missal genuinely loses a Sunday. The Liturgy of the Hours is organised by the same numbered weeks
and is not mentioned in the claim at all. And the Missal states the omission independently in
rubric~2\,b, which means the rule is not held by the Lectionary alone even as a matter of text.

\textbf{Where that leaves the claim.} The counterargument defeats any version of the claim that says
the omission is \emph{merely} a Lectionary event, and this reference does not say that. What
\loc{nualc}~44's second paragraph does is extend an existing reckoning --- \latin{eadem ratione},
``by the same reckoning'' --- to the other books. It presupposes the reckoning; it does not state it.
Neither the Norms nor the Missal anywhere supplies the criterion by which the omitted week is
identified, and neither supplies the reason. Both are supplied once, in \loc{olm}~104.3 with its
footnotes. So the claim stands in this narrowed form: the omission binds all three books, and the
Norms say so; but a reader who needs to know which week is omitted, and why, has exactly one place to
go, and it is the Lectionary. The unequal practical weight follows independently from rubric~3\,b and
is not affected by the counterargument at all.

\textbf{What would change it.} Any of the following, and this reference would revise:

\begin{itemize}
\item the \latin{Institutio generalis de Liturgia Horarum} stating the omission rule with a criterion
of its own --- this reference has not read that Instruction, which is a real and acknowledged gap in
the argument's evidence base rather than a rhetorical concession;
\item an authentic interpretation, a decree of the competent dicastery, or a later typical edition of
the Missal restating the rule with its criterion and reason;
\item evidence about the 1969 first typical edition of the \latin{Ordo lectionum Missae}. The 1981
decree of promulgation, protocol CD 240/81 of 21 January 1981, states as the first of the second
edition's five differences from the first that \latin{Textus ``Praenotandorum'' auctus est} --- the
text of the Praenotanda has been enlarged. Whether nn.~103--104 and their footnotes stood in 1969 or
entered in 1981 is therefore genuinely open, and the 1969 edition was not examined here; or
\item a general permission, in force in some territory, to join the readings of an omitted week to an
adjacent week, which would dissolve the practical asymmetry even if the textual argument survived.
\end{itemize}
\end{synthesis}

\begin{synthesis}{the leap-year dates in the \latin{Tabella temporaria}}

\textbf{Claim.} The one-day error in the February dates of the five leap-year rows reported in
Section~\ref{sec:tabella} more probably originated at or before the printed typical edition than in
the digital reproduction through which it was read --- but this is a judgment of probability and is
not established.

\textbf{Reasons.} The defect is perfectly correlated with a single arithmetic condition: a date
earlier than 29 February in a leap year. Six leap years fall in the table's range; the five whose Ash
Wednesday precedes the intercalated day are wrong in both February columns, and the sixth, whose Ash
Wednesday falls on 8 March, is right in every column. No transcription accident produces that
pattern; a date routine that adds a fixed day offset to a year-length assumption produces it exactly.
A table of this kind is itself the output of such a routine. Moreover, the reproduction consulted
here has visibly \emph{lost} the table's column structure, rendering it as running lines: that is the
signature of extraction and re-flowing from a set source, and extraction preserves values while
destroying layout. A reproducer who had regenerated the table computationally would have had no
reason to emit it as prose.

\textbf{The strongest counterargument, at full strength.} The artifact is a secondary digital
reproduction from a non-ecclesiastical host, of unknown production history, which is already
demonstrably incomplete --- its first appendix carries a heading and no contents --- and which has
demonstrably altered at least one element of the printed presentation. An object that has silently
dropped an entire appendix and destroyed a table's layout has forfeited any presumption of fidelity
in that table's values. It is entirely possible that its producer rebuilt the chronological table
from a script, precisely because the original ruled table was hard to reproduce in flowing text, and
that the script carried the leap-year bug. On that account the defect belongs to the reproduction,
the printed \latin{editio typica} is innocent, and this reference would be publishing a false
imputation against a promulgated book on the strength of a bootleg PDF.

\textbf{Where that leaves the claim.} The counterargument is strong enough that the claim is stated
here as a probability and not as a finding, and nothing in this reference depends on it. What does
not depend on it is the operational conclusion, which holds on either account: a February date read
from that table in a leap year must be recomputed, and the table's week numbers --- which are correct
in every row, including the defective ones --- may be relied on.

\textbf{What would change it.} Page images of the corresponding leaf of a licensed 2002
\latin{editio typica tertia}, or of the 2008 \latin{reimpressio emendata}, or the corresponding table
in an approved territorial edition of the Missal. Any one of those would settle the question in a
single reading. A published erratum or a \emph{Notitiae} correction touching the
\latin{Tabella temporaria} would settle it equally.
\end{synthesis}
